\section{Observed Trajectory Graphs and Exact Reuse}

Let an observed agent run be represented as a trajectory graph
\begin{equation}
T = (V,E,A),
\end{equation}
where each node $v \in V$ is an executed operation, each edge $e \in E$ records a dependency or call relation, and $A$ is the set of run artifacts such as events, reports, cache metadata, usage, errors, and redactions.

For a first implementation, the observed node kinds can be deliberately small:
\begin{equation}
\tau(v) \in \{\textsc{AgentStep}, \textsc{ModelCall}, \textsc{ToolCall}, \textsc{Handoff}\}.
\end{equation}
Each node carries a stable input hash, an optional output hash, timestamps, status, and metadata. Model and tool nodes additionally carry cache keys. This is the shape of the current \texttt{MigakiGraph}: \texttt{agent\_step}, \texttt{model\_call}, \texttt{tool\_call}, and \texttt{handoff} nodes connected by \texttt{depends\_on}, \texttt{calls}, and \texttt{produces} edges.

\begin{definition}[Operation cache key]
For an operation node $v$, define a cache key
\begin{equation}
K(v)=(op(v), name(v), H_{in}(v), H_{dep}(v), H_{run}(v)),
\end{equation}
where $op$ is the operation kind, $name$ is the model or tool name, $H_{in}$ is the normalized input hash, $H_{dep}$ is the dependency hash, and $H_{run}$ is the runtime hash containing implementation, SDK, model parameters, tool version, and output-schema assumptions.
\end{definition}

Two nodes from different runs are exact reuse candidates when their cache keys match and their prior output is admissible:
\begin{equation}
\Reusable(v_i, v_j)=1 \iff
K(v_i)=K(v_j) \wedge
status(v_i)=\textsc{Ok} \wedge
\Accept(output(v_i), v_j)=1.
\end{equation}

The acceptance predicate is intentionally explicit. For deterministic tools, it may be a contract check plus dependency hash equality. For model calls, it may require an output schema, validator result, replay policy, and privacy policy. If those requirements are missing, the node may be marked cacheable or potentially reusable, but it should not be replayed automatically.

This framing separates three claims that are often blurred:

\begin{itemize}[leftmargin=1.5em]
  \item \textbf{Observed:} a trajectory graph records what happened.
  \item \textbf{Reusable:} two nodes match under exact cache-key and dependency criteria.
  \item \textbf{Avoided:} a later run safely replayed a prior result instead of executing the node.
\end{itemize}

Only the third claim justifies cost or latency savings. The first two are already valuable: they turn agent execution from an opaque sequence of calls into an inspectable graph of reusable work.
